\chapter{Concluzii}
\label{cap:concluzii}

\section{Concluzii generale}

Lucrarea a pornit de la o întrebare pragmatică — se poate construi, din componente accesibile, un robot capabil să susțină activități terapeutice reale pentru copii neurodivergenți? — iar răspunsul demonstrat practic este afirmativ. TriSense există, funcționează cap-coadă pe hard\-ware real și acoperă un ciclu complet de interacțiune: copilul apasă un buton, vorbește, robotul îl înțelege, îi răspunde cu voce și mișcare, îi propune jocuri, îi verifică vizual construcțiile LEGO și consemnează totul pentru terapeut.

Din perspectivă tehnică, principala concluzie este că arhitectura distribuită pe trei niveluri — fiecare dispozitiv făcând strict ceea ce știe cel mai bine — a fost decizia care a făcut restul posibil: hub-ul LEGO nu știe că există internet, ESP32 nu știe ce spune modelul de limbaj, PC-ul nu atinge niciun motor. Această separare a făcut fiecare componentă testabilă izolat și a localizat imediat fiecare defect, de la redarea sacadată (rezolvată prin buffering) până la interferența modelului generativ cu logica jocului (rezolvată prin izolarea stărilor). Din perspectivă aplicativă, sistemul transpune fidel principiile din literatură: predictibilitatea și repetabilitatea care fac copilul disponibil pentru învățare \cite{scassellati2012, dautenhahn2004}, materialul LEGO intrinsec motivant \cite{legoff2004, owens2008}, feedback-ul imediat și exclusiv pozitiv, și rolul de mediator, nu de substitut. Merită consemnată și o concluzie metodologică: puterea modelelor generative vine la pachet cu nevoia de a le îngrădi explicit — prin prompturi atent formulate, praguri calibrate empiric și protejarea fluxului de control al jocurilor de creativitatea modelului conversațional.

\section{Direcții de dezvoltare ulterioară}

Câteva direcții se conturează natural: \textbf{validarea clinică extinsă} — primele sesiuni reale cu un copil neurodivergent au decurs încurajator, dar confirmarea eficacității cere un studiu amplu, cu protocol formal și avize etice, în colaborare cu terapeuți specializați în LBT; \textbf{completarea tripticului funcțiilor executive} cu jocurile aflate în stadiu de prototip (N-back verbal și Stroop adaptat); \textbf{detecția stării emoționale din voce} (analiza prozodiei, pentru a sesiza frustrarea înainte de escaladare); \textbf{înlocuirea ferestrei fixe de înregistrare} cu detecția automată a activității vocale; \textbf{funcționarea offline} prin modele locale; și o \textbf{interfață web pentru terapeut} (configurarea sesiunii, urmărirea în timp real, istoricul de metrici per copil).

Închei cu observația care mi se pare cea mai valoroasă din întregul proiect: niciuna dintre componentele folosite — setul LEGO, plăcile ESP32, serviciile de inteligență artificială — nu este exotică sau scumpă. Tot ce desparte un set educațional de un instrument terapeutic este integrarea. Iar integrarea, spre deosebire de hard\-ware-ul specializat, se poate replica oriunde există un calculator, o rețea Wi-Fi și cineva dispus să o facă.
